% ============================================================
% CHAPTER 5: LEARNING GOALS
% MCDA 5585 / 5586 Major Project Report
% Bhavik Kantilal Bhagat | A00494758 | Saint Mary's University
% ============================================================

\chapter{Learning Goals}
\label{chap:learning_goals}

The following learning goals were established at the start of my internship in
consultation with my academic and industry supervisors. They are organized into
technical and non-technical categories to reflect the dual competency development
expected of an MCDA industry placement.

% ────────────────────────────────────────────────────────────
\section{Internship Parameters}
% ────────────────────────────────────────────────────────────

My internship spanned \textbf{March~15 -- August~31, 2026}, for a total of
\textbf{900 hours} of logged work. The engagement operated in two phases:
a part-time phase from March~15 to April~15 (approximately 10--20 hours per
week) during which I completed the second semester of coursework in
parallel, followed by a full-time phase from April~15 to August~31
(approximately 40+ hours per week) upon academic semester completion. A
Mitacs BSI renewal (September~1 -- December~31, 2026) was subsequently
approved, extending the industrial collaboration into a second phase not
covered by this primary submission.

I was embedded within the \textbf{Intelligent Automation and
Site Reliability Engineering (IA-SRE)} team at CTM Halifax, operating
across three interlocking work streams:

\begin{itemize}
    \item \textbf{SRE \& Operational Observability} --- establishing SLIs
          and SLOs, applying the four golden signals to financial services
          workloads, and building a centralized observability platform on
          Amazon Managed Grafana (AMG).

    \item \textbf{Intelligent Automation \& Predictive Operations} ---
          authoring CloudFormation IaC templates, engineering a
          Lambda-based dynamic alarm provisioner, and developing the
          \texttt{ia-sre-resources-inventory} REST API (IISS-507).

    \item \textbf{Advanced Monitoring \& Capacity Analytics} ---
          integrating Container Insights for ECS telemetry, designing
          CloudWatch Logs Insights queries, and contributing to RPA
          production support rotations.
\end{itemize}

% ────────────────────────────────────────────────────────────
\section{Technical Learning Goals}
% ────────────────────────────────────────────────────────────

\begin{itemize}
    \item \textbf{AWS Cloud Services \& Infrastructure.}
          Gain production-level proficiency with the AWS services central to
          Citco's infrastructure: Amazon CloudWatch (metrics, alarms, Logs
          Insights, Metric Insights SQL); AWS Lambda (serverless functions,
          memory tuning, cold-start analysis, X-Ray tracing); Amazon ECS
          with Container Insights and Amazon EKS for containerized
          workloads; Amazon EC2 with CloudWatch Agent instrumentation;
          Amazon MSK (managed Kafka, consumer lag monitoring); Amazon SQS
          (queue-depth telemetry); Amazon SNS (alert routing); Amazon S3;
          AWS CloudFormation and AWS SAM; and AWS Secrets Manager.

          A deliberate learning goal was exposure to the \textbf{full
          spectrum of database categories} used in production financial
          platforms: Amazon DocumentDB; Amazon Neptune (graph database,
          CitcoWorks Knowledge Graph); Amazon RDS / Aurora; Amazon
          ElastiCache for Redis / Valkey (Meridian ODL tier); and Amazon
          OpenSearch Service. Additional infrastructure explored includes
          Amazon ECR, AWS CodeCommit, AWS CodePipeline, and the AWS IAM
          permission model (including the 16-role agent permission structure
          required for Grafana cross-account access).

    \item \textbf{Site Reliability Engineering Principles.}
          Operationalize SRE concepts in a financial services context ---
          defining SLIs and SLOs for RPA automation workloads, applying the
          four golden signals (latency, traffic, errors, saturation) to
          heterogeneous infrastructure, and designing proactive alerting
          thresholds that reduce mean time to detection (MTTD) for
          fund-impacting incidents.

    \item \textbf{Infrastructure-as-Code.}
          Author parameterized CloudFormation YAML templates for CloudWatch
          alarm deployment, develop a Lambda-based dynamic provisioner for
          per-resource alarm creation, and navigate the AWS Secrets Manager
          64\,KB payload limit encountered during provisioning (resolved by
          migrating the template payload to S3).

    \item \textbf{Observability Platforms.}
          Construct and operate a production-grade Amazon Managed Grafana
          (AMG) workspace --- including upgrading from version 10.4 to 12.4
          mid-delivery --- integrating CloudWatch as a data source, building
          14 dashboard rows with 110+ panels, and enabling Container
          Insights on five of six Meridian ECS clusters.

    \item \textbf{Proactive vs.\ Reactive Reliability Engineering.}
          Build proactive reliability systems: set SLO-aligned alarm
          thresholds (MSK consumer lag $< 10{,}000$ messages, ECS CPU
          $< 80\%$, Lambda error rate $< 1\%$ over 24h), and design the
          auto-healing pipeline (IISS-488) where alarm triggers initiate
          automated recovery actions rather than only notifications.

    \item \textbf{Python Development.}
          Write production-quality Python using \texttt{boto3} for AWS API
          interaction, \texttt{pytest} with Hypothesis for property-based
          testing (562-test suite, 99\% coverage), \texttt{Flask} with
          async route handlers, and \texttt{structlog} for structured JSON
          logging.

    \item \textbf{Software Engineering Practices.}
          Practice test-driven development (TDD) and spec-driven development
          via the Kiro IDE workflow (\texttt{.kiro/specs/} with
          \texttt{requirements.md}, \texttt{design.md}, \texttt{tasks.md}),
          enabling agentic sub-agent execution for code generation and PR
          creation.

    \item \textbf{AI-Augmented Engineering.}
          Leverage the Kiro IDE (Claude Sonnet), Amazon~Q for contextual
          code suggestions, and spec-driven agentic workflows to accelerate
          design documentation, code scaffolding, and automated test
          generation.
\end{itemize}

% ────────────────────────────────────────────────────────────
\section{Non-Technical Learning Goals}
% ────────────────────────────────────────────────────────────

\begin{itemize}
    \item \textbf{Agile and Kanban Methodology.}
          Work within Citco's JIRA-based delivery process: sprint planning,
          ticket decomposition, story-point estimation, and work-log
          tracking across the IISS, CAIS, and other project boards.

    \item \textbf{Technical Communication.}
          Produce engineering-grade written deliverables: incident RCA
          reports, architecture design documents, Confluence knowledge-base
          pages, and this academic report.

    \item \textbf{Production Incident Response.}
          Participate in four on-call RPA support rotations (Weeks 8, 12,
          16, 20) and handle time-sensitive production incidents including
          the CAIS Deadline hotfix (IISS-706), the MESO month-end stuck
          task retrigger (IISS-741, Critical), and VPL pricing failures
          (IISS-753, IISS-876).

    \item \textbf{Cross-Functional Collaboration.}
          Build working relationships across Halifax IA-SRE, Innovation IT
          Hyderabad, business operations teams, and the RPA support tier.
          Coordinate across Halifax EST and Hyderabad IST time zones.

    \item \textbf{Financial Domain Knowledge.}
          Develop working familiarity with fund administration operations:
          VPL queue (5~PM EST deadline), Corporate Actions workflows, MESO
          month-end cycles, and CAIS regulatory filing deadlines.
\end{itemize}
